%{% raw %}
\DocumentMetadata
  {
    lang=en-US,
    pdfversion=2.0,
    pdfstandard=ua-2,
    tagging=on,
  }

\documentclass{article}

\title{Theorem examples with order key}

\newtheorem{thm}{Theorem}[section]

\begin{document}

\section{A section}

\DebugTemplatesOn  % to see the processing in the log

\begin{thm}[It has a name]
  A theorem with default layout.
\end{thm}

\EditInstance{thmstyle}{plain}{
    order         =  {title,separator-a,number,punct,separator-b,note}
  , caption-decls = \sffamily\bfseries
  , note-decls    = \mdseries
  , after-hspace  = 20pt plus 2pt
}  

\begin{thm}[It has a name]
  A theorem with period after number and not after the note and a wide
  space separating it from th body text.
\end{thm}

\EditInstance{thmstyle}{plain}{
  , order       =  {number,separator-a,title,separator-b,note,separator}
  , separator-b = \\
  , separator   = \\[5pt]
  , after-hspace = 0pt
  , note-format = #1
  , note-decls = \footnotesize
}  

\begin{thm}[It has a name too]
  A theorem with number and title swapped and the note on a line by
  itself (if present).  No claim that this looks good :-)
\end{thm}

\begin{thm}
  As before, but without a note. Therefore \texttt{separator-b} in
  front of the non-existent note is dropped but the final
  \texttt{separator} remains.
\end{thm}

\end{document}

%{% endraw %}
